% chapters/abstract-cn.tex - 中文摘要

\chapter*{摘要}
\addcontentsline{toc}{chapter}{摘要}

\textbf{目的：}在规范化储备重复次数（RIR）自我指导已具备科学基础与可及性的前提下，基于计算机视觉的移动端数智化监控方案是否具备值得投入额外成本的适应性增量与实施可行性，仍缺乏针对性实证证据。本研究通过工具效度验证、预试验性随机对照训练比较与半结构化访谈三阶段递进设计，系统评估自研移动端计算机视觉系统 SportSci Pro 在青年男性下肢抗阻训练数智化监控中的测量学基础、干预方案可行性、探索性干预效果与用户接受度，并为后续正式随机对照试验提供效应量估计与流程优化参数。

\textbf{方法：}研究一招募 13 名有训练经验的健康青年男性（与研究二样本相互独立），在受控递增负荷实验条件下同步采集 SportSci Pro 与 GymAware 线性位置传感器的杠铃后蹲平均向心速度数据（60 对配对），采用组内相关系数（ICC）、Bland-Altman 分析与平均绝对误差（MAE）评估测量一致性。研究二采用两组平行、开放标签、分层区组随机对照预试验设计，将 36 名完成随机化的健康青年男性训练者（AI 辅助组 18 人，自我指导组 18 人）纳入 4 周、每周 2 次的下肢抗阻训练干预。

\textbf{结果：}研究一 ICC = 0.940（MAE = 0.047 m/s）；研究二真实场景 Rep 级配对 ICC = 0.991。PP 样本训练完成率 98.4\%，AI 组速度目标命中率 80.2\%，两组总重复次数差 $-$5.19 次（$p$ = 0.035）。CMJ 调整后差值 +2.16 cm（$g$ = 0.85），训练自我效能差值 +9.88 分（$p$ $<$ 0.001）。

\textbf{结论：}SportSci Pro 为低成本客观速度反馈方案下沉大众训练场景提供了测量学基础；4 周干预未产生确证性组间优效，但 CMJ 上捕获到方向性信号，训练自我效能显著提升。方案应定位为具有应用潜力、仍需大样本长周期验证的预试验性训练支持工具。

\textbf{关键词：}数智化监控；抗阻训练；基于速度的训练；预试验性研究；混合方法研究；随机对照预实验
